\documentclass{article}
\usepackage{graphicx}
\usepackage[margin=1.5cm]{geometry}
\usepackage{amsmath}
\usepackage{enumitem}

\setlist{nosep,leftmargin=*}

\begin{document}
\twocolumn

\title{In-Class Exercise: Python f-Strings and Unit Conversions}
\author{Prof. Jordan C. Hanson}
\date{}

\maketitle

\section{Formatted Strings}

\begin{enumerate}

\item Python \textbf{f-strings} allow numerical values to be inserted
into strings and control how those values are displayed.

\begin{itemize}

\item Try the following statements in the interpreter:

\begin{verbatim}
x = 3.14159265
print(f"{x:f}")
print(f"{x:.4f}")
print(f"{x:8.4f}")
print(f"{x:08.4f}")
\end{verbatim}

The format \texttt{f} uses fixed-point notation,
\texttt{.4f} displays four digits after the decimal,
\texttt{8.4f} uses a field at least eight characters
wide, and \texttt{08.4f} fills unused space with zeros.

\item Integer values can be formatted in a similar way:

\begin{verbatim}
n = 12345678
print(f"{n:d}")
print(f"{n:,d}")
print(f"{n:10d}")
print(f"{n:010d}")
\end{verbatim}

Predict the output of each statement before running it.

\end{itemize}
\end{enumerate}

\section{In-Class Exercise}

\begin{enumerate}

\item Write a Python program called
\texttt{unit\_conversions.py} that calculates and prints a table of British-to-metric unit conversions. Use 1 ft. = 0.3048 m, 1 mile = 1.609344 km, 1 lb. = 0.45359237 kg, 1 stone = 14 lb., 1 kg = 1000 gm, and 1 mph = 1.609344 km/hr.  Perform the following conversions:

\begin{center}
\begin{tabular}{c|c}
\textbf{Input} & \textbf{Convert to} \\ \hline
125 ft & meters \\
3.75 miles & kilometers \\
175 lb & kilograms \\
12 stone & kilograms \\
2.5 lb & grams \\
0.5 stone & grams \\
65 mph & kilometers/hour
\end{tabular}
\end{center}

\begin{itemize}

\item Store each starting value in a variable and calculate each
converted value using Python arithmetic.

\item Distance and speed results must contain \textbf{3 significant
figures}.  For these particular values, use the appropriate
fixed-point precision:
\begin{verbatim}
38.1
6.04
105
\end{verbatim}

Thus, different rows may require \texttt{.1f}, \texttt{.2f},
or \texttt{.0f}.

\item Weight results must contain \textbf{4 significant figures}.
For these particular values, the results should have forms such as

\begin{verbatim}
79.38
76.20
1134
3175
\end{verbatim}

Again, select the appropriate number of decimal places.

\item Use field widths so that the numerical results line up in
columns.  For example,

\begin{verbatim}
print(f"Feet to meters    "
      f"{feet:8.1f}"
      f"{meters:10.1f}")
\end{verbatim}

\item Your final output should have the general form

\begin{verbatim}
Unit Conversion Table

Conversion          Input     Result
Feet to meters      125.0       38.1
Miles to km          3.75       6.04
Lb to kg            175.0      79.38
Stone to kg          12.0      76.20
Lb to grams           2.50       1134
Stone to grams        0.50       3175
Mph to km/hr          65.0        105
\end{verbatim}

Do not type the converted results directly into the
\texttt{print()} statements.  Python must calculate each result
before the f-string formats it.

\end{itemize}
\end{enumerate}

\end{document}